\origpage{2}{}
\textbf{1.2} \textit{If $L$ is a vital $5$-linkage in a simple graph $G$, then either some two terminals in different components of $L$ are adjacent in $G$, or some terminal has the same degree in $G$ and in $L$, or $G$ is null.}

\fix{We give a proof of (1.2). Suppose that $G$ is not null and that the first two outcomes do not hold. Then each terminal $t$ of $L$ has larger degree in $G$ than in $L$, and so there is an edge $e \notin E(L)$ of $G$ incident with $t$; let its other end be $u$. Then $u \ne t$, since $G$ is simple.}

\fix{\hspace*{12pt}(1) \textit{For each such $t$ and $u$, $u$ is an internal vertex of a path of $L$ that does not contain $t$.}}

\fix{\textit{Subproof.} Suppose that $u$ belongs to the component $P$ of $L$ that contains $t$. Then $P$ has an edge, since $u \ne t$, and $t$ is an end of $P$; let $P$ have vertices $t = v_0, v_1, \dots, v_m$ in order, and let $u = v_j$ where $j \ge 1$. If $j = 1$ then $e$ and the edge of $P$ between $v_0$ and $v_1$ are parallel edges of $G$, which is impossible. If $j \ge 2$, let $P'$ be obtained from $P$ by replacing its subpath between $v_0$ and $v_j$ by $e$. Then $P'$ is a path with the same ends as $P$, and so replacing $P$ by $P'$ gives a linkage in $G$ with the same pattern as $L$, and different from $L$ since it contains $e$. This contradicts that $L$ is vital. Hence $u \notin V(P)$. Since $V(L) = V(G)$, $u$ belongs to another component $Q$ of $L$, and $u$ is not a terminal of $L$ since the first outcome does not hold. Hence $u$ is an internal vertex of the path $Q$. This proves (1).}

\fix{A path of $L$ that has an internal vertex has two distinct ends, both terminals of $L$. Since $L$ has at most five terminals, it follows that}

\fix{\hspace*{12pt}(2) \textit{There do not exist three vertices of $L$, none of them a terminal of $L$, that belong to three different components of $L$.}}

\fix{Since $G$ is not null, some vertex of $G$ is not a terminal of $L$ (if $v \in V(G)$ is a terminal, apply (1) to $v$). Let $y$ be such a vertex, and let $P$ be the component of $L$ that contains it; then $P$ is a path with distinct ends $a, b$. By (1), $a$ and $b$ are adjacent in $G$, by edges not in $E(L)$, to internal vertices $x_1, y_1$ of paths of $L$ different from $P$. By (2) applied to $y, x_1, y_1$, the vertices $x_1$ and $y_1$ belong to the same component $Q$ of $L$. Let $Q$ have ends $c, d$. By (1) again, $c$ and $d$ are adjacent in $G$, by edges not in $E(L)$, to internal vertices $x_2, y_2$ of paths of $L$ different from $Q$; and by (2) applied to $y, x_1, x_2$ and to $y, x_1, y_2$, both $x_2$ and $y_2$ belong to $P$. Let $P'$ be the path formed by the edge between $a$ and $x_1$, the subpath of $Q$ between $x_1$ and $y_1$, and the edge between $y_1$ and $b$; and let $Q'$ be the path formed in the same way from $c$, the subpath of $P$ between $x_2$ and $y_2$, and $d$. These are paths, since $a, b \notin V(Q)$ and $c, d \notin V(P)$. They are disjoint, since $P'$ has no vertex outside $V(Q) \setminus \{c, d\}$ except $a, b$, and $Q'$ has no vertex outside $V(P) \setminus \{a, b\}$ except $c, d$. Hence replacing $P, Q$ by $P', Q'$ gives a linkage in $G$ with the same pattern as $L$, different from $L$ since the edge between $a$ and $x_1$ is not in $E(L)$. This contradicts that $L$ is vital, and so proves (1.2).}\fixnote{fix:1.2-proof-added}{The paper states (1.2) without proof. The argument above is the one used in the formalization.}

This raises the question of whether (1.1) is true in general with tree-width replaced by path-width; and indeed it is, as we shall show in section 12.

Our proof of (1.1) is as follows. From a theorem of [8] it follows immediately that every graph with a vital $p$-linkage has no $K_n$ minor, where $n \ge \frac{5}{2}p + 1$. Consequently we can apply the results of [9, 10] concerning the structure of graphs excluding a fixed minor. They ``almost'' tell us that $G$ has bounded genus. The remainder of the proof falls into two main parts; first we prove it when $G$ really does have bounded genus, and then we fix the gaps implied by ``almost``.

\sectionhead{2 Some basic lemmas}

In this section we establish some lemmas about vital linkages that we shall need repeatedly. We use $\setminus$ to denote the result of deletion; thus, $G \setminus X$ is the graph obtained from $G$ by deleting $X$.

\textbf{2.1} \textit{If $L$ is a vital $p$-linkage in $G$, and $X \subseteq V(G)$, then $L \setminus X$ is a vital $(p + 2|X|)$-linkage in $G \setminus X$.}

The proof is clear.

A \textit{separation} of $G$ is a pair $(A, B)$ of subgraphs of $G$ with union $G$ and with $E(A \cap B) = \emptyset$; its \textit{order} is $|V(A \cap B)|$.

\textbf{2.2} \textit{If $L$ is a linkage in $G$ with set of terminals $X$, and $(A, B)$ is a separation of $G$, then $L \cap B$ is a linkage in $B$ with set of terminals a subset of $(X \setminus V(A)) \cup V(A \cap B)$.}

Again, the proof is clear, as is the proof of the next lemma.

\textbf{2.3} \textit{If $L$ is a linkage in $G$, and $(A, B)$ is a separation of $G$, and $L'$ is a linkage in $B$ with the same pattern as $L \cap B$, then $(L \cap A) \cup L'$\fix{, with the vertices that are internal vertices of paths of $L \cap B$ deleted from $L \cap A$,} is a linkage in $G$ with the same pattern as $L$.\fixnote{fix:2.3}{As printed, the first claim fails when a vertex of $V(A \cap B)$ is internal in $L \cap B$ and $L'$ does not use it. That vertex is then a one-vertex component of $L \cap A$, and it stays one in the union. For example, let $G$ be the 4-cycle $a, v, b, x$, let $L$ be the path $a, v, b$, let $A$ have vertices $a, v, b$ and no edges, let $B = G$, and let $L'$ be the path $a, x, b$. The inserted deletion removes no edge, since a vertex that is internal in $L \cap B$ has degree $0$ in $L \cap A$. The later uses of (2.3) only need a linkage inside $(L \cap A) \cup L'$ with the pattern of $L$. The second claim holds as printed.} In particular, if $L$ is a vital linkage in $G$ then $L \cap B$ is a vital linkage in $B$.}

\fix{We give the proof of (2.3). Let $L_A$ be obtained from $L \cap A$ by deleting the vertices that are internal vertices of paths of $L \cap B$, and let $M = L_A \cup L'$. Since $(A, B)$ is a separation, every edge of $L$ belongs to exactly one of $L \cap A$, $L \cap B$; so a vertex that is internal in $L \cap B$ has degree $0$ in $L \cap A$, and $L_A$ has all the edges of $L \cap A$. Since $L'$ and $L \cap B$ have the same pattern, they have the same terminals, and each terminal has the same degree in both (degree $0$ if its block has one member, and $1$ otherwise). In particular, every vertex of $V(L_A) \cap V(B)$ is a terminal of $L \cap B$ and of $L'$, since it is in $V(L \cap B)$ and is not internal there.}

\fix{\hspace*{12pt}(1) \textit{Every vertex of $M$ has degree $\le 2$ in $M$; every vertex of $M$ that is in $V(L_A)$ or is a terminal of $L \cap B$ has the same degree in $M$ and in $L$; and $M$ and $L$ have the same terminals.}}

\fix{\textit{Subproof.} Let $v \in V(M)$. If $v$ is in $V(L_A)$ or is a terminal of $L \cap B$, then its edges in $L_A$ are its edges in $L \cap A$, and its edges in $L'$ are as many as its edges in $L \cap B$, by the remarks above (both numbers are $0$ if $v \notin V(L \cap B)$, because then $v \notin V(L')$, as every vertex of $V(L') \cap V(L_A)$ and every terminal of $L'$ is in $L \cap B$). So its degree in $M$ is its degree in $L$. Otherwise $v \in V(L')$ is not a terminal of $L'$, and has degree $2$ in $L'$ and in $M$. This proves the first two statements. For the third, a terminal $t$ of $L$ is either in $V(L \cap A)$ and not internal in $L \cap B$, so in $V(L_A)$, or in $V(L \cap B)$ with degree $\le 1$, so a terminal of $L \cap B$; and in both cases $t$ has degree $\le 1$ in $M$. Conversely, a vertex of degree $\le 1$ in $M$ is in $V(L_A)$ or is a terminal of $L'$, that is, of $L \cap B$, and so has the same degree in $L$. This proves (1).}

\fix{\hspace*{12pt}(2) \textit{Let $H \subseteq M$ have at least one edge, such that for each component $C$ of $L'$, every vertex of degree $1$ in $H \cap C$ is an end of $C$. Then there exists $H_L \subseteq L$ with at least one edge, such that a vertex has degree $1$ in $H_L$ if and only if it has degree $1$ in $H$.}}

\fix{\textit{Subproof.} Let $C$ be a component of $L'$ with an edge in $H$. The components of $H \cap C$ with an edge are subpaths of $C$ whose two ends have degree $1$ in $H \cap C$, and so are ends of $C$; hence $H \cap C = C$. Let $H_L$ be the union of $H \cap L_A$ and, for each component $C$ of $L'$ with $C \subseteq H$, the path of $L \cap B$ with the same ends as $C$ (it exists since $L'$ and $L \cap B$ have the same pattern). Then $H_L \subseteq L$, and $H_L$ has an edge, since $H$ has one. An end of such a path has one edge in it and one edge in $C$, and its other edges in $H_L$ and in $H$ are those of $H \cap L_A$. An internal vertex of such a path has degree $2$ in $H_L$; it is not in $L_A$, and not an end of a component of $L'$, so it does not have degree $1$ in $H$. Every other vertex of $H$ or $H_L$ has the same edges in $H_L$ as in $H \cap L_A$, and in $H$ it has in addition either no edge of $L'$ or two (if it is internal in some $C \subseteq H$, and then it is not in $L_A$, by the last remark of the first paragraph). This proves (2).}

\fix{Now $M$ is a linkage in $G$. For by (1) it has maximum degree $\le 2$, and if it had a circuit $H$, then (2) would apply to $H$: a vertex of degree $1$ in $H \cap C$ has an edge of $H$ in $L_A$, and so is a terminal of $L'$, that is, an end of $C$. Then $H_L \subseteq L$ has an edge and no vertex of degree $1$, which is impossible. By (1), $M$ and $L$ have the same terminals. Let $u \ne w$ be terminals in the same component of $M$, and let $R$ be the path of $M$ between them. By the same argument (2) applies to $R$, since $u, w$ are terminals of $L$, and so are terminals of $L'$ if they belong to $L'$. So there exists $H_L \subseteq L$ in which $u, w$ are the only vertices of degree $1$; since $L$ has maximum degree $\le 2$, the component of $H_L$ that contains $u$ is a path, and its other end is $w$. Hence $u, w$ are in the same component of $L$. Since a terminal of degree $0$ in $M$ has degree $0$ in $L$ by (1), every block of the pattern of $M$ is a block of the pattern of $L$, and since the two patterns partition the same set, they are equal.}

\fix{For the second statement, suppose that $L$ is vital. Then $V(L \cap B) = V(B)$, and $L \cap B$ is a linkage in $B$. Let $L'$ be a linkage in $B$ with the same pattern as $L \cap B$. By the first statement, $M$ is a linkage in $G$ with the same pattern as $L$, and so $M = L$. An edge of $L \cap B$ is therefore in $L_A$ or in $L'$, and not in $L_A$ since $E(A \cap B) = \emptyset$; so $E(L \cap B) \subseteq E(L')$, and $E(L') \subseteq E(M) \cap E(B) = E(L \cap B)$. A vertex of $L \cap B$ is either a terminal of $L \cap B$, and so of $L'$, or an internal vertex, which is in $V(M)$ but not in $V(L_A)$, and so is in $V(L')$; and $V(L') \subseteq V(M) \cap V(B) = V(L \cap B)$. Hence $L' = L \cap B$, and so $L \cap B$ is vital.}\fixnote{fix:2.3-proof-added}{The paper says that the proof of (2.3) is clear. The argument above is the one used in the formalization.}

We need the operation of ``splitting'' a vertex of a graph. (For the following to make sense, we regard a graph as a triple consisting of a set of vertices, a set of edges, and an appropriate incidence relation between them.) Let $v$ be a vertex of a graph $G$, and let $\delta_1, \delta_2 \subseteq E(G)$ with $\delta_1 \cap \delta_2 = \emptyset$, so
